\section{Introduction}\label{sec:intro}
A tool-using agent usually learns that something went wrong from the tool itself: a timeout, an
error code, a malformed reply. The harder case is a tool that reports success but did something
else. A write is dropped, lands on the wrong field, or is applied only in part, and the result the
agent sees is exactly what a correct call would have returned. Agents often fail to detect such quiet
failures~\cite{kraishan2026loud}, benchmarks find implicit tool failures harder than explicit
ones~\cite{zhu2026tools}, what an agent's history says about a step can hurt when the world does not
match it~\cite{pakhomov2026cleaning}, and recent systems add monitors, contracts or verification layers
to catch such failures~\cite{panthi2026outcome,babu2026selfhealing,dip2026invariantaware,kaynar2026tools}.

These systems bundle two decisions: \emph{when} to suspect a step and \emph{what} to do about it.
When they are compared, a better trigger and a better recovery procedure cannot be told apart. We
separate them. We hold the agent, its prompt and its recovery message fixed, and vary only the trigger
that decides when the agent is told that a step may be wrong. We compare seven triggers, from
``never'' to an independent read-back of the world state after every write, on matched tasks where
every arm draws the same fault and the same fault slot.

This is a narrow controlled comparison on a mock record service, not a new method. Runtime monitors of
action postconditions already trigger replanning in classical planning~\cite{cardoso2026rvreplan},
and verification of non-atomic tool calls~\cite{mansoor2026verified}, handling of writes that may
already have taken effect~\cite{li2026where}, outcome contracts~\cite{panthi2026outcome} and state
checkers for wrong reads~\cite{zhang2026rite} exist for LLM agents. Outcome Monitors already vary when
a fixed generic warning attaches, find that the recovery-tool list carries the gain, and detect no
difference from timing~\cite{panthi2026outcome}. Our study differs in its faults, silent writes graded on the final
world state; in including an independent world-state read-back alone and combined with a return-value
check; in the full trigger set under recover, with abstain for the two checking triggers; and in a
pre-registered analysis on a small local model, with a larger hosted model as an exploratory contrast. Our contributions are:
\begin{itemize}
  \item A matched, pre-registered comparison of seven triggers and two responses across seven fault
        types, graded on the final world state (Sections~\ref{sec:method} and~\ref{sec:setup}).
  \item On Qwen3-8B, read-back beats a return-value contract on silent write faults by \RPOneDiff{} in
        success (Holm $p$ \RPOneHolmP{}), and combining both beats never recovering by
        \RPThreeDiff{} (Holm $p$ \RPThreeHolmP{}) (Section~\ref{sec:results}).
  \item An inconclusive pre-registered test: the contract does not significantly raise success on
        wrong reads (Holm $p$ \RPTwoHolmP{}).
  \item Exploratory evidence that detection and repair are separate problems: read-back flags silent
        write faults on all three models, but only the strongest repairs most of them.
\end{itemize}
